\documentclass[a4paper,twoside,11pt]{article}

\usepackage[spanish]{babel}
\usepackage[utf8]{inputenc}
\usepackage{graphicx}
\usepackage{float}
\usepackage[skins]{tcolorbox}
\usepackage{titlepic}
\usepackage{fancyhdr}
\usepackage{booktabs}
\usepackage[margin=2.5cm]{geometry}
\usepackage[hidelinks]{hyperref}

\pagestyle{fancy}
\lhead{Geotecnia para geólogos}
\rhead{\thepage}
\cfoot{Programa}
\renewcommand{\headrulewidth}{0.4pt}
\renewcommand{\footrulewidth}{0.4pt}

\graphicspath{{assets/img/}}

\title{GEOTECNIA PARA GEÓLOGOS}
\author{Prof.: Edier Aristizábal\\[5ex]
\includegraphics[width=8.0cm]{unal2}}
\date{}

%################################BODY############################################
\begin{document}
\maketitle

\begin{center}
\emph{versión del programa}: \today
\end{center}

\section*{Introducción}
El curso de Geotecnia está orientado a estudiantes de Ingeniería Geológica que deseen comprender el
comportamiento mecánico e hidráulico de los suelos y de los macizos rocosos, y su aplicación en las obras
de ingeniería y en la evaluación de amenazas. La geotecnia nació como una rama de la ingeniería civil, pero
sus objetos de estudio son materiales geológicos: suelos transportados y residuales, saprolitos, macizos
fracturados y depósitos de vertiente. El curso busca formar geólogos capaces de traducir una descripción
geológica en un modelo geotécnico, con parámetros, mecanismos de falla y factores de seguridad, y de
dialogar con los ingenieros sobre los supuestos que hay detrás de un diseño.

El curso se estructura en cuatro bloques: (i) los materiales geológicos (suelos, rocas y agua),
(ii) esfuerzos y deformaciones (estado de esfuerzos, esfuerzos efectivos, distribución de esfuerzos,
compactación, consolidación y cimentaciones), (iii) resistencia al corte de suelos y rocas, y
(iv) laderas y estructuras de contención (movimientos en masa, estabilidad de taludes y empujes de tierras).
Los ensayos geotécnicos de campo y de laboratorio se estudian a lo largo del curso, a medida que se
abordan los temas relacionados.

\section{PROGRAMA}
El contenido del curso es:

\section*{Parte I: Los materiales geológicos}

\subsection{Introducción a la geotecnia}
Definición. Historia de la geotecnia. Obras típicas. El papel del geólogo en el estudio geotécnico.
Título H de la NSR-10.

\subsection{Mecánica de suelos}
Meteorización. Suelos residuales y transportados. Perfil de meteorización. Catena. Minerales de arcilla.
Relaciones volumétricas y gravimétricas. Granulometría de suelos gruesos y finos. Límites de Atterberg.
Clasificación de suelos para ingeniería (USCS). Suelos colapsables y expansivos.

\subsection{Mecánica de rocas}
Roca intacta y macizo rocoso. Discontinuidades y su descripción (ISRM). Clasificaciones geomecánicas:
RQD, RMR, Q y GSI. Modelo geotécnico de macizos rocosos.

\subsection{El agua en el suelo y en el macizo rocoso}
Ciclo hidrológico. Estados del agua en el suelo. Tensión superficial y capilaridad. Suelos parcialmente
saturados y succión. Curva característica. Flujo unidimensional y ley de Darcy. Conductividad hidráulica.
Flujo bidimensional y redes de flujo. Fuerzas de filtración. Diseño de filtros. Flujo en medio poroso y en
medio fracturado.

\section*{Parte II: Esfuerzos y deformaciones}

\subsection{Estado de esfuerzos}
Definiciones. Tensor de esfuerzos. Análisis bidimensional. Círculo de Mohr y método del polo.
Diagramas p-q y trayectorias de esfuerzos.

\subsection{Esfuerzos efectivos}
Resistencia de materiales. Curva esfuerzo-deformación. Principio de esfuerzo efectivo. Perfiles de
esfuerzos. Flujo vertical y gradiente crítico.

\subsection{Distribución de esfuerzos verticales}
Principio de superposición. Esfuerzo por carga puntual, lineal, de franja, circular y rectangular.
Bulbo de presiones. Carta de Newmark.

\subsection{Compactación}
Métodos de mejoramiento de suelos. Ensayos Proctor. Curva de compactación. Control en campo.

\subsection{Consolidación y asentamientos}
Compactación frente a consolidación. Ensayo edométrico. Curva de compresibilidad. Preconsolidación.
Cálculo de asentamientos. Teoría de la consolidación unidimensional.

\subsection{Cimentaciones}
Tipos de cimentaciones. Capacidad de carga. Modos de falla. Teoría de Terzaghi y ecuación general.
Factor de seguridad. Capacidad de pilotes. Asentamientos inmediatos.

\section*{Parte III: Resistencia}

\subsection{Resistencia al cortante de los suelos}
Criterios de falla. Ecuación de Coulomb. Criterio de Mohr-Coulomb. Cohesión y ángulo de fricción.
Dilatancia. Resistencia pico y residual. Ensayo de corte directo. Ensayos triaxiales CD, CU y UU.
Parámetros de Skempton. Compresión inconfinada.

\subsection{Resistencia de rocas}
Movilización de la resistencia por cohesión y fricción. Factores que influyen en la resistencia.
Falla progresiva. Criterios de falla de Griffith, Mohr-Coulomb y Hoek-Brown. Resistencia de
discontinuidades (Barton-Bandis).

\section*{Parte IV: Laderas y estructuras de contención}

\subsection{Movimientos en masa}
Definiciones y nomenclatura. Partes, geometría y volumen. Clasificación de Cruden y Varnes (1996).
Movimientos en masa detonados por lluvia. Otras clasificaciones. Cartografía.

\subsection{Estabilidad de laderas y taludes}
Fuerzas estabilizadoras y desestabilizadoras. Factor de seguridad. Talud infinito con y sin
infiltración. Talud finito con superficie plana y circular. Método de dovelas: ordinario, Bishop
simplificado, Janbu y métodos rigurosos. Ábacos. Análisis cinemático de macizos rocosos. Análisis
probabilístico del factor de seguridad: métodos FOSM y de estimativas puntuales.

\subsection{Empujes de tierras y muros de contención}
Tipos de muros. Estados en reposo, activo y pasivo. Teorías de Rankine y Coulomb. Efecto del agua y de
las sobrecargas. Estabilidad de muros de contención.

\subsection*{Ensayos geotécnicos}
Este tema se cubre a lo largo del curso por parte de los estudiantes, a medida que se abordan los temas
relacionados con cada ensayo, y forma parte de la participación en clase: toma de muestras, SPT, CPT,
ensayos de permeabilidad (Lefranc y Lugeon), prospección geofísica, ensayos in situ, ensayos de
laboratorio (límites, compresión simple, corte directo, triaxiales, consolidación) e instrumentación.

\section{EVALUACIÓN DEL CURSO}
\begin{center}
\begin{tabular}{p{8.5cm}cp{4.5cm}}
\toprule
Componente & Peso & Momento \\
\midrule
Trabajo 1: perfil de meteorización de una ladera (informe escrito) & 10\% & Al terminar la teoría de esfuerzos \\
Trabajo 1: presentación & 20\% & Al terminar la teoría de esfuerzos \\
Trabajo 2: análisis de estabilidad de la misma ladera (informe escrito) & 10\% & Al terminar estabilidad de laderas \\
Trabajo 2: presentación & 20\% & Al terminar estabilidad de laderas \\
Parcial teórico sobre todo el contenido del curso & 30\% & Al terminar toda la teoría \\
Participación en clase & 10\% & A lo largo del curso \\
\bottomrule
\end{tabular}
\end{center}

\textbf{Trabajo 1 (30\%).} Cada grupo selecciona una ladera y describe su perfil de meteorización: unidades
geológicas, horizontes según Dearman o Deere y Patton, propiedades índice y clasificación de los materiales,
condiciones de agua y perfil de esfuerzos totales, presiones de poros y esfuerzos efectivos. Se entrega un
informe escrito (10\%) y se presenta ante el grupo (20\%) luego de ver la teoría de esfuerzos.

\textbf{Trabajo 2 (30\%).} Sobre la misma ladera, se construye el modelo geológico-geotécnico y se realiza un
análisis de estabilidad para diferentes escenarios: condición seca, diferentes posiciones del nivel freático
o saturación por lluvia, y sismo (análisis pseudoestático), solos y combinados. Se entrega un informe escrito
(10\%) y se presenta ante el grupo (20\%).

\textbf{Parcial (30\%).} Examen teórico sobre todo el contenido del curso, que se realiza al terminar toda
la teoría.

\textbf{Participación (10\%).} Participación en clase, consultas y ejercicios.

\section*{Horario y aula}
Horario: por definir. \quad Aula: por definir.

\section{MATERIAL DEL CURSO}
\begin{itemize}
\item Página del curso: \url{https://edieraristizabal.github.io/Geotecnia/}
\item Libro del curso: \url{https://edieraristizabal.github.io/Geotecnia/libro/}
\item Diapositivas: \url{https://edieraristizabal.github.io/Geotecnia/html/Geotecnia.html}
\end{itemize}

\section{REFERENCIAS}
El curso utilizará material de diferentes fuentes bibliográficas, entre las cuales se destacan las
siguientes, por lo cual se recomienda su consulta:

\begin{itemize}
\item Fundamentos de ingeniería geotécnica. Braja M. Das. Thomson Learning. 2001.
\item Ingeniería geológica. Luis I. González de Vallejo. Pearson. 2002.
\item Mecánica de suelos. Juárez Badillo y Rico Rodríguez. Limusa. 2005.
\item Mecánica de suelos. Peter L. Berry y David Reid. McGraw-Hill. 1993.
\item Soil mechanics in engineering practice. Karl Terzaghi, Ralph Peck y Gholamreza Mesri. John Wiley \& Sons. 1996.
\item Practical rock engineering. Evert Hoek. RocScience. 2007.
\item Soil strength and slope stability. J. Michael Duncan y Stephen G. Wright. John Wiley \& Sons. 2005.
\item Slope stability and stabilization methods. L. Abramson, T. S. Lee, S. Sharma y G. M. Boyce. John Wiley \& Sons. 2002.
\item Landslides: investigation and mitigation. Special Report 247. Transportation Research Board. 1996.
\item Movimientos en masa en la región andina: una guía para la evaluación de amenazas. Proyecto Multinacional Andino (GEMMA). 2007.
\item Reglamento Colombiano de Construcción Sismo Resistente NSR-10, Título H. 2010.
\end{itemize}

\end{document}
